\documentclass[12pt]{article}
\usepackage{amsmath, amssymb, geometry, setspace}
\geometry{margin=1in}
\setstretch{1.2}
\setlength{\parindent}{0pt}


\title{Practice Problem Set 3}
\author{ECON 320 -- Introduction to Mathematical Economics}
\date{Fall 2025}

\begin{document}
\maketitle



\section*{1. Nonlinear Production and Cost Optimization}

A firm produces output according to:
\[
Q(L) = 50L^{0.6}e^{-0.02L}
\]
where \(L\) is labour.

\begin{enumerate}
\item Determine the value of \(L^*\) that maximizes \(Q(L)\).
\item Verify that the second-order condition holds.
\item Interpret economically what happens to output beyond \(L^*\).
\item Compute the elasticity of \(Q\) with respect to \(L\) at the optimum.
\end{enumerate}



\section*{2. Consumer Utility with Exponential Preferences}

A consumer’s utility function is given by:
\[
U(x) = 5x^{0.4}e^{-0.1x}
\]

\begin{enumerate}
\item Find \(x^*\) that maximizes utility.
\item Show that utility first increases and then decreases with \(x\).
\item Compute the maximum utility value \(U(x^*)\).
\item Explain intuitively why the exponential term reflects satiation.
\end{enumerate}



\section*{3. Logarithmic Profit Maximization}

A firm’s profit function is:
\[
\pi(Q) = 200\ln(Q + 1) - 30Q
\]

\begin{enumerate}
\item Derive the first-order condition for maximization.
\item Solve for the optimal output \(Q^*\).
\item Verify that the second-order condition holds.
\item Compute the profit at \(Q^*\).
\item Interpret how logarithmic revenue affects marginal profit.
\end{enumerate}



\section*{4. Advertising Response with Diminishing Returns}

Sales as a function of advertising \(A\) are:
\[
S(A) = 400(1 - e^{-0.05A}) - 2A
\]

\begin{enumerate}
\item Derive the optimal advertising level \(A^*\).
\item Check that the solution satisfies the second-order condition.
\item Compute the corresponding maximum sales \(S(A^*)\).
\item Discuss the meaning of \(e^{-0.05A}\) in capturing saturation effects.
\end{enumerate}



\section*{5. Elasticity via Logarithmic Differentiation}

The cost of production is given by:
\[
C(Q) = 10Q^{0.7}(2 + \ln Q)
\]

\begin{enumerate}
\item Use logarithmic differentiation to find the elasticity of cost with respect to output.
\item Evaluate the elasticity when \(Q = 10\).
\item Interpret the result: is cost increasing at an increasing or decreasing rate?
\end{enumerate}



\section*{6. Maximization with Rational Function}

The firm’s revenue is modeled by:
\[
R(P) = \frac{100P}{4 + P^2}
\]
where \(P\) is the price per unit.

\begin{enumerate}
\item Find the price \(P^*\) that maximizes total revenue.
\item Verify that it is indeed a maximum.
\item Determine the corresponding revenue \(R(P^*)\).
\item Explain the trade-off between price and quantity implied by this function.
\end{enumerate}



\section*{7. Cost Minimization with Logarithmic Form}

A firm’s cost function is:
\[
C(x) = 20x + 50\ln(x)
\]
for \(x>0\).

\begin{enumerate}
\item Find the value of \(x^*\) that minimizes cost.
\item Confirm that it is a minimum.
\item Interpret the economic meaning of the logarithmic term in cost.
\item Compute the elasticity of cost at \(x^*\).
\end{enumerate}



\section*{8. Growth Model with Continuous Time}

A population grows according to:
\[
P(t) = 800e^{0.03t} - 400e^{-0.02t}
\]

\begin{enumerate}
\item Find the time \(t^*\) at which the population grows fastest.
\item Compute the maximum growth rate.
\item Interpret the combination of exponential growth and decay components.
\item Explain why \(t^*\) corresponds to a turning point in \(P(t)\).
\end{enumerate}



\section*{9. Logarithmic Demand and Revenue}

The demand curve is:
\[
P(Q) = 60 - 5\ln(Q + 1)
\]
and total revenue is \(R(Q) = P(Q)Q\).

\begin{enumerate}
\item Derive the first-order condition for maximizing revenue.
\item Solve for the revenue-maximizing quantity \(Q^*\).
\item Verify that \(R''(Q^*) < 0\).
\item Discuss the role of logarithmic demand in moderating price sensitivity.
\end{enumerate}



\section*{10. Challenge: Nonlinear Utility and Elasticity}

A utility function is given by:
\[
U(x) = \ln(3x + 2) - 0.5x^2
\]

\begin{enumerate}
\item Determine the value of \(x^*\) that maximizes utility.
\item Check that \(U''(x^*) < 0\).
\item Compute the elasticity of utility with respect to \(x\) at the optimum.
\item Provide an economic interpretation: why does marginal utility turn negative for large \(x\)?
\end{enumerate}



\end{document}
